\documentclass[11pt,letterpaper]{article}
\usepackage[margin=1in]{geometry}
\usepackage[T1]{fontenc}
\usepackage{amsmath,amsthm,booktabs,array,longtable,graphicx}
\usepackage{newtxtext,newtxmath}
\usepackage[round,authoryear]{natbib}
\usepackage{setspace,xr-hyper,xurl}
\usepackage[hidelinks]{hyperref}
\usepackage{fancyhdr}
\onehalfspacing\setlength{\emergencystretch}{2em}
\newtheorem{theorem}{Theorem}[section]
\newtheorem{proposition}[theorem]{Proposition}
\newtheorem{lemma}[theorem]{Lemma}
\newtheorem{corollary}[theorem]{Corollary}
\theoremstyle{definition}\newtheorem{remark}[theorem]{Remark}
\newcommand{\R}{\mathbb R}\newcommand{\E}{\mathbb E}
\newcommand{\argmax}{\operatorname*{arg\,max}}
\newcommand{\argmin}{\operatorname*{arg\,min}}
\pagestyle{fancy}\fancyhf{}\fancyhead[L]{\small Finite-Catalog Resource Allocation}
\fancyhead[R]{\small Anonymous}\fancyfoot[C]{\thepage}\setlength{\headheight}{14pt}

% BEGIN retained/current source: revisions/or-r52-resource-path-20260925/generated/metrics.tex
\newcommand{\FineMaxGap}{0.000500}
\newcommand{\FineMinTime}{3.523}
\newcommand{\FineMaxTime}{19.762}
\newcommand{\FineMaxRelative}{1.63\%}
\newcommand{\FineMinPacked}{0.216}
\newcommand{\FineMaxPacked}{6.628}
\newcommand{\FineMaxDenBits}{553}
\newcommand{\AllResourceExact}{31}
\newcommand{\GroupNOneComplete}{2}
\newcommand{\MIPMaxSeconds}{0.038}
\newcommand{\MIPMaxGap}{$5.55\,10^{-17}$}

% END source


% BEGIN retained/current source: revisions/or-r54-exact-price-path-20260926/generated/metrics.tex
\newcommand{\PriceExact}{10}
\newcommand{\PriceTolerance}{19}
\newcommand{\PriceLimited}{1}
\newcommand{\PriceChecked}{30}
\newcommand{\PriceMinBytes}{541}
\newcommand{\PriceMaxBytes}{5895}
\newcommand{\AnytimeExact}{10}
\newcommand{\ForcedAnchors}{45}
\newcommand{\HardnessSubsets}{1692}
\newcommand{\HardnessBooks}{3004}

% END source

\hypersetup{pdftitle={Finite-Catalog Resource Allocation: A Tariff-Sensitive Complexity Frontier},pdfauthor={Anonymous}}
\externaldocument{.build/r60/electronic_companion}[electronic_companion.pdf]
\begin{document}\hypersetup{pageanchor=false}
\begin{titlepage}\centering{\Large\bfseries Finite-Catalog Resource Allocation: A Tariff-Sensitive Complexity Frontier\par}
\vspace{0.25in}Anonymous manuscript for Operations Research --- R60\par\vspace{0.2in}
\begin{minipage}{\textwidth}\textbf{Abstract.} A service provider must allocate a fixed aggregate promise while using a small shared catalog of terminal commands. Each history has an expected cap and a ceiling on every realized obligation, and each installed command incurs a fixed charge. We derive an exact, implementable resource-path representation and establish a tariff-sensitive complexity frontier. Exact optimization is parameterized-hard in the command allowance and cap count, even with linear rewards and free service. Under those same physical primitives, a capacity dynamic program is fixed-parameter tractable in the number of commands whose opening fees differ from a standard tariff. Uniform fees therefore permit polynomial-time optimization without a resource grid. We also characterize conservation on length-bounded paths, retain an original-feasible additive approximation, and identify when a price bound supports one implementable book. A prospectively specified multi-budget study separates rational certification from numerical solver bounds and retains every interrupted request. The analysis shows how tariff structure, rather than the command allowance alone, can determine tractability, while exact representations connect optimal values to feasible operational decisions.
\par\vspace{0.15in}\textbf{Keywords:} resource allocation; dynamic programming; parameterized complexity.
\par\vspace{0.1in}\textbf{Subject classifications:} Programming: resource allocation; Dynamic programming: finite-catalog design.
\par\vspace{0.1in}\textbf{Area of review:} Optimization.\end{minipage}\vfill September 27, 2026\end{titlepage}
\hypersetup{pageanchor=true}
\section{Introduction}
A finite command catalog creates a distinction between the resources allocated to a decision and the number of commands used to implement it. A provider may deliver different expected obligations to different customer histories while installing only a small common set of terminal commands. We call the installed set a command book. Each installed command incurs a fixed charge. Individual expected caps and ceilings on every realized obligation restrict how an accepted aggregate promise can be implemented. Choosing expected targets first and a catalog afterward can therefore miss both a discrete cost and a feasibility constraint.

This paper establishes a tariff-sensitive complexity frontier for joint resource allocation and implementation. The main contrast holds even when rewards are linear and preliminary service is free. With arbitrary command-specific opening fees, exact optimization remains parameterized-hard in both the command allowance and the number of distinct expected caps. With a standard fee and only a few exceptional commands, however, we obtain an exact fixed-parameter algorithm whose polynomial exponent does not depend on the number of exceptions. Uniform fees give a polynomial-time algorithm with no discretization of capacities or the accepted promise. Thus a small command allowance does not by itself make the problem tractable, whereas a simple installation tariff can do so without restricting the physical heterogeneity.

The positive result follows from a capacity representation inside the original model. A selected book's gross value depends only on its attainable aggregate terminal capacity. This capacity telescopes along consecutive selected commands. Conditional on the installed exceptional commands and the book size, the total opening charge is fixed; it is then sufficient to retain one maximum-capacity label at each state. A recovered label is not merely a numerical value: it supplies terminal lotteries and preliminary service that meet the original promise and every individual constraint. The same tables compute the provider's exact decision frontier as a uniform installation fee changes. The dynamic-programming principle is familiar; its applicability here follows from the model-specific capacity and fixed-charge identities.

The exact selected-boundary representation also supports the general concave model. We characterize capacity conservation on acyclic graphs and extend the recognition criterion to paths with an explicit edge allowance. Complementing resources merges feasible starting points into a single deficit state. Rounding and repair preserve the selected path, its command count and the accepted promise. These results retain the original-feasible additive guarantee even where the tariff algorithm is inapplicable. A corridor-rehabilitation model illustrates the conservation and repair mechanism outside renewal commands, with its physical quantities and separability assumptions stated explicitly in the companion.

Price-based certificates and resource approximation play different roles. We characterize exactly when one price-active book supports the accepted promise, and bound the residual gap by the distance to that book's maximizing target interval. This explains why price methods can be strong on root-supported instances without implying that their general search trees are small. A uniform positive installation fee can still create a price gap even though the tariff algorithm solves the problem exactly. The earlier deficit implementation's unfavorable timing results remain part of the evidence; an accuracy guarantee does not establish practical speed superiority.

The computational study evaluates these distinctions rather than claim operational adoption. It preserves the earlier records, distinguishes distinct model inputs from tolerance-tagged cases, and adds prospectively frozen standard-tariff, exception-tariff and curved-return families under three end-to-end allowances. The direct mixed-integer convex formulation has no service-envelope approximation. Numerical bounds from its solver remain separate from independently checked rational certificates. Setup, serialization and verification consume the same allowance as optimization; unsuccessful requests remain in the denominator.

\subsection{Relationship to prior work}
Lagrangian constrained-path bounds, label search and resource-state approximation are established tools \citep{HandlerZang1980,BeasleyChristofides1989,Fisher1981,Hassin1992,LorenzRaz2001,PuglieseGuerriero2013}. Contemporary bucket-graph labeling organizes dominance and reduced-cost elimination for routing subproblems \citep{Sadykov2021}; heuristic-guided constrained search addresses large networks and multiple resource limits \citep{Ahmadi2025}. Those methods are not displaced by our deficit dynamic program. Our graph edges instead carry allocatable continuous resources, and our claims concern an exact representation of individual obligations, a recognizable conservation restriction, and recovery of an original equality. Generic path algorithms do not establish those model identities. No timing comparison against a routing implementation is inferred from our allocation experiments.

Continuous allocation on a fixed book belongs to separable resource allocation \citep{Patriksson2008}. Fixed-charge nonlinear allocation also admits perspective reformulations and branch-and-Benders-cut approaches \citep{Yang2025}. These are important formulation and algorithmic antecedents, not results that opening fees generally become easy. Our positive theorem exploits the additional ordered-capacity identity and the fact that conditioning on exceptional commands makes a given book size have a common charge. Its hardness counterpart concerns parameterized complexity, rather than the presence of binary variables alone. Parameterized knapsack classifications distinguish structural parameters, numerical values and their encodings \citep{Gurski2019}; their bounds do not transfer to our shared-book constraints without a reduction. We supply that reduction directly inside the original policy model.

Small distributional support is not a small shared installation catalog. Moment-set extreme-point results bound individual supports \citep{Winkler1988}; they do not charge for their union. Auction menu complexity additionally imposes incentive constraints absent here \citep{HartNisan2013,Babaioff2016}. Choice-based assortment optimization changes demand through the offered set \citep{Rusmevichientong2010}, whereas our history probabilities are fixed and terminal probabilities are controlled. Our reward perturbation certificate holds feasibility inputs fixed; it is not a statistical distributional-robustness claim \citep{EsfahaniKuhn2018}.

Finally, concave convolution uses classical matrix-searching structure \citep{Aggarwal1987}; the finite-support proof permits holes created by merged sources. General medoid selection and one-dimensional weighted-median clustering are antecedent algorithms, not independent novelty claims \citep{Arya2004,Wu1991}. Their role is to make the model's oscillation-based grouping bound operational. Table~\ref{tab:frontier60} separates these antecedents from the results established here. The paper's scientific contribution is the exact representation and tariff-sensitive frontier, with conservation, approximation and certification supplying implementable consequences.
\section{Finite-Catalog Renewal Design}\label{sec:model52}
There are $k$ histories. History $j$ has probability $\pi_j>0$, expected-obligation cap $b_j\in[0,1]$, and realization ceiling $\tau_j\geq b_j$, with $\sum_j\pi_j=1$. Write $\bar b=\sum_j\pi_jb_j$. Intermediate service has cost $h_j:[0,b_j]\to\mathbb R_+$, continuous, convex and nondecreasing, with $h_j(0)=0$. The initial common terminal reward $f:[0,1]\to\mathbb R$ is strictly increasing and concave. The catalog is $\mathcal A=\{a_1<\cdots<a_N\}\subset[0,1]$. Opening command $a_i$ costs $\rho_i\geq0$.

A nonempty selected book $c=(c_1<\cdots<c_s)$ has $s\leq m$. After observing $j$, the provider selects deterministic pre-draw service $y_j\geq0$ and a lottery $C_j$ supported on $c$. The problem is
\begin{align}
 V^*(B)&=\max_{c,y,C}\left\{\sum_j\pi_j\bigl[\mathbb E f(C_j)-h_j(y_j)\bigr]-\sum_{z\in c}\rho(z)\right\},\label{eq:original52}\\
 t_j&=y_j+\mathbb E C_j\leq b_j,\qquad y_j+C_j\leq\tau_j\quad\text{almost surely},\qquad \sum_j\pi_jt_j=B.\label{eq:constraints52}
\end{align}
Here $B\in[0,\bar b]$ is accepted, rather than optimized away. The common-risk specialization is $\tau_j=b_j+\delta$. Expected and realized obligations are different restrictions. The terminal symbol does not carry a free history index, inherited service tier, or shared random seed. The cardinality $m$ measures the command alphabet, not total implementation storage: targets, encoder logic, probability precision, program constants, and certificates are additional resources. The mathematical results below concern the stated finite catalog; continuous-design results have a separate approximation error.

For a feasible book, let $d_j=\max(c\cap[0,\tau_j])$, and let $F_{j,c}$ interpolate $f$ linearly on this eligible prefix. A singleton prefix has the constant value $f(d_j)$ at its only point.
\begin{proposition}[Canonical response]\label{prop:canonical52}
A book is feasible if and only if $c_1\leq\min_jb_j$ and $c_1\leq B\leq\bar b$. At target $c_1\leq t\leq b_j$, its optimal response is
\begin{equation}\label{eq:canonical52}
 v_{j,c}(t)=F_{j,c}(\min\{t,d_j\})-h_j((t-d_j)_+).
\end{equation}
It is attained by $y_j=(t-d_j)_+$ and an adjacent eligible lottery with mean $\min(t,d_j)$. The response is concave in $t$.
\end{proposition}
\begin{proof}
Every feasible command satisfies $C_j\leq\tau_j$ since $y_j\geq0$. At mean $\mu=\mathbb E C_j$, concavity on the ordered support bounds the reward by $F_{j,c}(\mu)$. The payoff is at most $F_{j,c}(\mu)-h_j(t-\mu)$ for $c_1\leq\mu\leq\min(t,d_j)$. Both terms are nondecreasing in $\mu$, so the upper endpoint maximizes this bound. When $t\leq d_j$, adjacent interpolation has zero intermediate service and support below $\tau_j$. When $t>d_j$, use command $d_j$ surely and service $t-d_j$; the realized sum is $t\leq b_j\leq\tau_j$. Thus the bound is attained without weakening either restriction. Interpolation slopes decrease and are nonnegative; the tail slopes of $-h_j$ are nonpositive and nonincreasing. This proves concavity. Feasible targets fill the entire weighted interval $[c_1,\bar b]$, proving necessity and sufficiency.
\end{proof}

The algorithmic specialization is
\begin{equation}\label{eq:quadratic52}
 f(x)=rx-\tfrac12 qx^2,\quad r\geq q\geq0,\quad r>0,\qquad h_j(y)=\tfrac12\gamma_jy^2,\quad \gamma_j\geq0.
\end{equation}
All probabilities, caps, ceilings, catalog points, shape parameters, opening charges, and requested accuracies are rationally encoded in exact-arithmetic claims. Optimization identities allow real nonnegative charges; bit-complexity assertions do not assign a finite encoding to arbitrary real numbers. A common Lipschitz constant for the canonical responses is
\begin{equation}\label{eq:L52}L=\max\{r,\max_j\gamma_j\}>0.\end{equation}
Continuity of the cost on its compact domain makes the stated extrema attained; every rational quadratic instance already satisfies this regularity. A fixed-book allocation can be solved by filling positive chord slopes in decreasing order, then zero-cost service capacity, then quadratic service tails. The service breakpoints are $\gamma_j(b_j-d_j)_+$. This takes $O(ks+k\log(k+1))$ rational operations with sorted marginal events. Alternatively, a direct supporting-price sweep verifies
\begin{equation}\label{eq:fixeddual52}
 \sum_j\pi_jv_{j,c}(t_j)=\lambda B+\sum_j\pi_j\max_{c_1\leq t\leq b_j}\{v_{j,c}(t)-\lambda t\}.
\end{equation}
Eligible knots, caps, and clipped stationary tail points suffice to check each scalar maximum. This certificate proves optimality for that book, not over all books.
\subsection{Infeasible inputs and certificate scope}\label{sec:infeasible59}
For well-formed primitives and a nonempty catalog, feasibility is decided before an incumbent is initialized. If $m<1$, $B<0$, $B>\bar b$, or $a_1>\min\{B,\min_jb_j\}$, no feasible book exists. Necessity follows from nonemptiness, the aggregate caps and the minimum possible command; sufficiency follows by using the singleton $a_1$ and filling deterministic service to the accepted promise. An empty catalog is likewise infeasible. A malformed probability vector, unsorted catalog or invalid primitive is instead an input error and must not be reported as an optimization certificate.

An infeasibility record identifies the failed inequality and carries the external input digest. It has no incumbent, lower policy or finite objective interval. An independent checker recomputes the probability-weighted cap, minimum cap and minimum catalog point from the supplied original input. Validating a scalar inequality is sufficient; search is unnecessary. The current solver modules retain their legacy exception on infeasible requests, whereas \texttt{feasibility59.py} provides this explicit preflight record and checker. The new panel is generated feasible and does not count preflight rejection as an optimization success.
\paragraph{Units and loss.} Obligations and caps use one normalized resource unit. Rewards, service costs and installation charges share a net-value unit. An additive tolerance is stated in that value unit; rescaling every payoff and the tolerance by the same positive factor preserves its decision meaning. The present theory makes no empirical currency calibration.
\section{Constructive Selected-Boundary Decomposition}\label{sec:boundaries52}
This section assumes the common reward $f$ in Section~\ref{sec:model52}. For a selected book $c=(u_1<\cdots<u_s)$, write $a=u_1$ and $d_j=\max\{u_i:u_i\leq\tau_j\}$.
\begin{lemma}[Global terminal-first allocation]\label{lem:terminal52}
An optimal fixed-book canonical allocation cannot have both $t_i>d_i$ and $t_j<\min(b_j,d_j)$.
\end{lemma}
\begin{proof}
Choose $0<\delta\leq\min\{t_i-d_i,\pi_j[\min(b_j,d_j)-t_j]/\pi_i\}$. Reduce $t_i$ by $\delta$ and increase $t_j$ by $\pi_i\delta/\pi_j$. The weighted change is zero. The donor's terminal reward is unchanged and its service cost weakly decreases; the recipient's strictly increasing terminal interpolant improves. Both responses remain feasible by Proposition~\ref{prop:canonical52}, contradicting optimality.
\end{proof}
For selected neighbors $u<v$, define
\begin{align}
 \delta_j(u,v)&=\mathbf1_{\{\tau_j\geq v\}}[\min(b_j,v)-u]_+,
 &T_{uv}&=\sum_j\pi_j\delta_j(u,v),\label{eq:terminalarc52}\\
 J_{uv}&=\{j:u\leq\tau_j<v\},
 &Q_{uv}&=\sum_{j\in J_{uv}}\pi_j(b_j-u)_+.\label{eq:exit52}
\end{align}
The capped service value on an exit set is
\begin{equation}\label{eq:servicearc52}
 H_{uv}(y)=\min\left\{\sum_{j\in J_{uv}}\pi_jh_j(z_j):
 0\leq z_j\leq(b_j-u)_+,\ \sum_{j\in J_{uv}}\pi_jz_j=y\right\},\quad 0\leq y\leq Q_{uv}.
\end{equation}
Empty sets have zero capacity and zero value at zero. Let $\sigma_{uv}=[f(v)-f(u)]/(v-u)$ and
\begin{equation}\label{eq:arc52}
 \Phi_{uv}(x)=\sigma_{uv}\min(x,T_{uv})-H_{uv}((x-T_{uv})_+),
 \qquad 0\leq x\leq T_{uv}+Q_{uv}.
\end{equation}
For the terminal edge $(u,\dagger)$, use $T_{u\dagger}=0$, $J_{u\dagger}=\{j:\tau_j\geq u\}$, and $\Phi_{u\dagger}=-H_{u\dagger}$. The functions $\Phi$ are concave. Under the quadratic primitives, they are $L$-Lipschitz with $L$ as in \eqref{eq:L52}.

\begin{lemma}[Nested terminal layers]\label{lem:nested54}
Start each history at target $a$. Fill the selected terminal layers in increasing edge order. Once every preceding layer is full, each history with $\delta_j(u,v)>0$ is at target $u$. Every weighted amount $x\in[0,T_{uv}]$ in the next layer has a feasible individual allocation. After all terminal layers are full, history $j$ is at $\min(b_j,d_j)$; its remaining service capacity belongs to exactly one exit set.
\end{lemma}
\begin{proof}
A history with positive increment on $(u,v)$ has $b_j>u$ and $\tau_j\geq v$. It is eligible for all earlier selected commands, and the increments of those earlier layers telescope to $u-a$. This proves the first assertion, including a history whose cap will truncate the current layer.

For any desired $x$, initialize residual $r=x$ and visit histories in any fixed order. Give history $j$ increment $e_j=\min\{\delta_j(u,v),r/\pi_j\}$ and reduce $r$ by $\pi_je_j$. Since $x\leq\sum_j\pi_j\delta_j(u,v)$, the final residual is zero. Each affected target lies in $[u,\min(b_j,v)]$ and is implemented by the adjacent lottery on $u,v$ with upper probability $e_j/(v-u)$. It has no intermediate service and both realizations are below $\tau_j$. Histories with zero increments keep their previously feasible responses.

If a layer is partial, stop: its predecessors are full and its successors unused. If every layer is full, the increments for history $j$ telescope to $\min(b_j,d_j)-a$. It exits exactly after $d_j$, at an ordinary crossing edge or at the terminal edge. Thus its remaining capacity is $(b_j-d_j)_+$ and appears in exactly one $H$. Any feasible service vector on the exit sets implements target $\min(b_j,d_j)+z_j$. A positive $z_j$ requires $b_j>d_j$, and its implementation is the sure command $d_j$ plus service $z_j$; their sum is at most $b_j\leq\tau_j$. Zero-capacity histories require no service. The exit sets are disjoint, so these implementations combine without conflicting allocations.
\end{proof}

\begin{lemma}[Implementable rearrangement]\label{lem:rearrange54}
For a fixed selected path, any feasible arc-resource vector can be rearranged, with unchanged total resource and no smaller payoff, into full earlier terminal layers, at most one partial terminal layer, and service only after all terminal layers are full. That rearranged payoff is attained by original history-level policies.
\end{lemma}
\begin{proof}
Split the resource on each arc into its terminal and service parts in \eqref{eq:arc52}. If any terminal capacity remains unused while some service is positive, transfer the smaller of the available terminal capacity and a positive service amount to that terminal layer. The terminal gain is positive, while reducing service weakly lowers cost. Repeat until there is no such pair. This operation never changes charges or the sum of resources.

Next, if an earlier terminal layer is unfinished and a later one is positive, transfer the smaller of the unfinished capacity and the later mass. Selected chord slopes are positive and nonincreasing, so payoff cannot decrease. Filling layers successively yields the asserted form after finitely many exhausted donors or filled recipients. Each arc still respects its capacity: terminal mass is at most $T_e$, and service, when present, is at most $Q_e$ after that arc's terminal part is full. Allocate any remaining service by minimizers of the corresponding $H_e$; these minimizers exist by compactness and continuity. Lemma~\ref{lem:nested54} implements the resulting terminal pattern and the disjoint service allocations. Its gross payoff is precisely $f(a)+\sum_e\Phi_e(x_e)$ for the rearranged vector. An arbitrary original vector need not be implemented at its unrearranged payoff; the weak improvement is the comparison needed for exact optimal values.
\end{proof}

\begin{theorem}[Exact selected-boundary resource path]\label{thm:path52}
Let $P(c)$ contain the successive selected edges and the terminal edge of a feasible book. Then
\begin{align}
 \sum_{e\in P(c)}(T_e+Q_e)&=\bar b-a,\label{eq:capacityidentity52}\\
 V_c(B)&=f(a)-\sum_{z\in c}\rho(z)+
 \max_{\substack{0\leq x_e\leq T_e+Q_e\\\sum_e x_e=B-a}}\sum_{e\in P(c)}\Phi_e(x_e).\label{eq:pathidentity52}
\end{align}
Maximizing over paths with at most $m$ selected commands solves the original joint problem exactly. Reconstruction fills terminal layers in order, uses the greedy partial-layer allocation of Lemma~\ref{lem:nested54}, and then solves the disjoint capped service allocations. The resulting policy obeys the original promise, expected caps, realization ceilings, and command budget.
\end{theorem}
\begin{proof}
For history $j$, the terminal increments sum to $\min(b_j,d_j)-a$ and its unique service capacity is $(b_j-d_j)_+$. Their sum is $b_j-a$, proving \eqref{eq:capacityidentity52} after weighting. In an original optimum, Lemma~\ref{lem:terminal52} places service after terminal capacity. Among terminal layers, exchanging equal weighted amounts orders the common decreasing chord slopes; Lemma~\ref{lem:nested54} justifies the individual feasibility of that ordering. If all terminal capacity is full, the exit-set cost minimizations describe all remaining feasible service allocations. Hence an original optimum supplies an arc allocation of the same value. Conversely, Lemma~\ref{lem:rearrange54} implements a weakly improved version of every feasible arc allocation. The two comparisons yield equality of the attained maxima. Each selected charge is subtracted once, and the terminal edge does not open an additional command.
\end{proof}

\begin{corollary}[Endogenous eligibility]\label{cor:endogenous52}
A book of $s$ commands has at most $s$ nonempty selected-prefix classes. Inserting unselected candidates changes neither its arc functions nor its fixed-book optimum.
\end{corollary}
\begin{proof}
The nonempty exit sets partition histories according to their last eligible selected command. Their definitions and terminal increments depend on the selected endpoints, not intervening unused candidates.
\end{proof}
\subsection{An eligibility-independent capacity potential}
The path-capacity identity has a stronger local form. It explains why eligibility affects returns, but does not create a hidden resource state.
\begin{proposition}[Capacity potential]\label{prop:potential53}
Define $p(u)=\sum_j\pi_j\min(b_j,u)$ for catalog vertices and $p(\dagger)=\bar b$. Every selected-boundary arc satisfies
\begin{equation}\label{eq:potential53}
 T_{uv}+Q_{uv}=p(v)-p(u),\qquad Q_{u\dagger}=p(\dagger)-p(u).
\end{equation}
Thus its capacity is independent of the realization ceilings as long as $\tau_j\geq b_j$. For a feasible anchor $a$, $p(a)=a$.
\end{proposition}
\begin{proof}
Fix one history and an ordinary edge $u<v$. If $\tau_j<u$, then $b_j<u$, so both its contribution to the arc and $\min(b_j,v)-\min(b_j,u)$ are zero. If $u\leq\tau_j<v$, then $b_j<v$ and its exit-service capacity is $(b_j-u)_+=\min(b_j,v)-\min(b_j,u)$. If $\tau_j\geq v$, its terminal increment is $[\min(b_j,v)-u]_+$, which equals the same difference. These cases partition the histories. Multiplication by the probabilities and summation prove the first identity. For the terminal edge, $\tau_j<u$ again implies $(b_j-u)_+=0$, giving $Q_{u\dagger}=\sum_j\pi_j(b_j-u)_+=\bar b-p(u)$. Finally, feasibility implies $a\leq b_j$ for every history.
\end{proof}
Changing a ceiling can still change an arc's payoff through its exit set. Proposition~\ref{prop:potential53} asserts capacity invariance, not value continuity across a threshold. It also identifies the renewal reduction as a capacity-conservative configuration of the type used in the production model below.
\subsection{A heterogeneous-reward primal extension}\label{sec:packing54}
For heterogeneous rewards define the history-specific chord slope $\sigma^j_{uv}=[f_j(v)-f_j(u)]/(v-u)$. This gives the appropriate primal arc payoff. Simply replacing the common reward by different rewards in \eqref{eq:arc52} would be incorrect, because one chord slope no longer describes all histories. Instead, define the within-edge terminal allocation value
\begin{equation}\label{eq:terminalprofile54}
 R_{uv}(x)=\max\left\{\sum_{j:\tau_j\geq v}\pi_j\sigma^j_{uv}e_j:
 0\leq e_j\leq\delta_j(u,v),\ \sum_j\pi_je_j=x\right\},\quad 0\leq x\leq T_{uv}.
\end{equation}
The feasible within-edge polytope is nonempty and compact at each stated total, so its linear objective attains a maximum. Its value is continuous, increasing and concave, with marginal slopes drawn from the history-specific chord slopes. Set
\begin{equation}\label{eq:heteroarc54}
 \widetilde\Phi_{uv}(x)=R_{uv}(\min(x,T_{uv}))-H_{uv}((x-T_{uv})_+).
\end{equation}
The terminal edge still has payoff $-H_{u\dagger}$. All terminal slopes are positive and all service marginal returns are nonpositive, so this joined payoff is concave.

\begin{theorem}[History-level component packing]\label{thm:packing54}
For a common ordered catalog with ceiling-induced prefix eligibility and heterogeneous increasing concave terminal rewards, replace $f(a)$ in \eqref{eq:pathidentity52} by $\sum_j\pi_jf_j(a)$ and every $\Phi$ by $\widetilde\Phi$. The resulting scalar-resource path identity is exact. Its capacities remain those in \eqref{eq:capacityidentity52}. For any feasible arc-resource vector, within-edge optimizers followed by history-level packing produce a feasible original policy with the same individual component totals and no smaller payoff.
\end{theorem}
\begin{proof}
For a fixed book, introduce a component $e_{j,uv}\in[0,\delta_j(u,v)]$ for every eligible terminal segment, and service $z_j\in[0,(b_j-d_j)_+]$. Temporarily drop the requirement that a history's earlier components must be filled before its later ones. The relaxed objective, apart from charges, is
\[
 \sum_j\pi_j\left[f_j(a)+\sum_{(u,v)}\sigma^j_{uv}e_{j,uv}-h_j(z_j)\right],
 \qquad \sum_j\pi_j\left[\sum_{(u,v)}e_{j,uv}+z_j\right]=B-a.
\]
Every original canonical policy embeds in this relaxation by filling its eligible segments in order and then its service tail. Conversely, fix any relaxed component allocation and define
$t_j=a+\sum_{(u,v)}e_{j,uv}+z_j$.
The sum of that history's component capacities is exactly $b_j-a$, so $a\leq t_j\leq b_j$. Its components can be packed in their original order while preserving $t_j$: move service into any unfinished positive terminal segment, then move mass from a later positive segment into any unfinished earlier segment. Each move weakly improves payoff, because the history's terminal slopes are decreasing and positive and service cost is nondecreasing. This comparison is entirely within a history; it needs no common ordering across histories. At termination the components implement its canonical response at $t_j$. Proposition~\ref{prop:canonical52} supplies the adjacent lottery or sure-command service policy, with every realization constraint intact. The weighted promise is unchanged by packing. Thus relaxation and original problem have the same optimum.

Now regroup the relaxed components by selected edge, assigning $z_j$ to the unique exit edge of history $j$. Groups share only the scalar total-resource equality; their individual component bounds are otherwise disjoint. For a fixed group resource $x$, optimal allocation fills positive terminal returns first, whose optimum is $R_{uv}$, and then minimizes the exit-set service cost $H_{uv}$. This is exactly \eqref{eq:heteroarc54}. Maximizing over group resources is therefore the same relaxation already proved exact. The component-capacity telescoping is unchanged. Within-edge optimizers followed by the preceding history-level packing provide the asserted constructive recovery.
\end{proof}

In the quadratic specialization, $R_{uv}$ is evaluated by sorting the eligible history slopes and filling weighted segment capacities. Its Lipschitz constant is at most $\max_j r_j$. The service part is bounded by $\max_j\gamma_j$. Hence the polynomial additive resource construction also extends to heterogeneous rewards with these new kernels. The heterogeneous merged-origin implementation in Section~\ref{sec:deficit56} evaluates these primal kernels directly. The predecessor readers retain the earlier common-reward comparison and its separate scope.
\section{A Recognizable Class of Resource Paths}\label{sec:class59}
The renewal reduction is an instance of a more general allocation problem. Let $G=(V,E)$ be a finite directed acyclic graph with a common sink $t$, a set of possible starting vertices $S$, capacities $C_e\geq0$, continuous edge payoffs $\phi_e$ on $[0,C_e]$, and source payoffs $A_s$. A decision selects a source-to-sink path $P$ of at most $h\geq1$ edges and allocations $0\leq x_e\leq C_e$ satisfying $\sum_{e\in P}x_e=B_s$. Its value is $A_s+\sum_{e\in P}\phi_e(x_e)$, with any fixed edge or vertex charge included once in the corresponding path payoff. Retain only vertices reachable from a source and able to reach the sink. No ordered catalog, history, or eligibility prefix is required for this abstract problem.

\begin{theorem}[Recognition and source-independent deficits]\label{thm:recognition59}
The following are equivalent on the retained graph: (i) for every source, every source-to-sink path has the same total capacity; (ii) every vertex has a well-defined suffix capacity $H_v$, with $H_t=0$ and $C_{uv}=H_u-H_v$ on every edge. The condition can be recognized in $O(|V|+|E|)$ arithmetic operations. When it fails, two paths from a common source with different total capacities form a directly checkable witness.

Under these equivalent conditions, replacing $x_e$ by $w_e=C_e-x_e$ yields a single terminal resource level independent of source and path if and only if $H_s-B_s=R$ is the same constant for every admissible source. In that case the exact optimum is
\[
 \max_{s,P}\left\{A_s+\max_{0\leq w_e\leq C_e,\;\sum_{e\in P}w_e=R}
                  \sum_{e\in P}\phi_e(C_e-w_e)\right\}.
\]
All sources can be initialized at deficit zero. Infeasible sources with $B_s\notin[0,H_s]$, or with no path within the edge limit, are excluded explicitly; if none remain, the problem is infeasible.
\end{theorem}
\begin{proof}
Condition (ii) telescopes, proving (i). Conversely fix a vertex $v$ and a source-to-$v$ prefix. Concatenating that prefix with any two $v$-to-sink suffixes gives valid paths, because the graph is acyclic. Condition (i) makes their suffix capacities equal. Define their common sum as $H_v$. Appending any edge $uv$ to a suffix from $v$ yields $H_u=C_{uv}+H_v$.

In reverse topological order set $H_t=0$. At $u$, compute $C_{uv}+H_v$ for each successor and require all results to agree. This examines each edge once. Store a representative suffix at every vertex by successor pointers. If two values disagree, their edges followed by the stored suffixes have different capacities; prepend any retained source prefix. The resulting witness has two explicit paths, whose capacity sums can be checked independently. Reachability and the minimum number of edges to the sink are also computed by topological passes.

On any path starting at $s$, $\sum C_e=H_s$, so $\sum x_e=B_s$ is equivalent to $\sum w_e=H_s-B_s$. This proves necessity and sufficiency for this particular complement coordinate, and is a bijection between the constrained allocation boxes. Continuations at a vertex depend only on that vertex, elapsed edge count, and accumulated deficit. Keeping the best prefix at each such state therefore loses no feasible complete path or value. The predecessor chain retains the actual source. The statement is not a claim that no other algorithm can merge states when capacity conservation fails.
\end{proof}

\begin{theorem}[Additive approximation on conservative paths]\label{thm:abstractapprox59}
Suppose the recognized instance has $R\geq0$, every admissible source has $H_s\geq R$, and the centered deficit payoffs
$g_e(w)=\phi_e(C_e-w)+\lambda_0w$ are concave and $K$-Lipschitz for one common $K>0$. With $\eta=\varepsilon/(2Kh)$, a layered path dynamic program on deficit multiples of $\eta$, accepting totals in $[R-h\eta,R]$, returns a feasible policy and a value interval of width at most $\varepsilon$. Apart from kernel evaluation, its arithmetic count is
$O(h|E|Q\log(Q+1)+h|V|Q)$, where $Q=\lfloor R/\eta\rfloor+1$. The result preserves the source, path, edge limit, and the exact original equality. Without concavity, the same guarantee holds for Lipschitz kernels using the direct $O(h|E|Q^2)$ convolution.
\end{theorem}
\begin{proof}
Round the deficits of an optimal path down. At most $h$ coordinates change, their total change is less than $h\eta$, and their centered payoff loss is at most $Kh\eta$. The best accepted grid value $Z_\eta$ consequently gives upper bound $Z_\eta-\lambda_0R+Kh\eta$. For a maximizing grid path, increase deficits to total $R$. Its total available capacity is $H_s\geq R$, so a greedy fill within residual edge capacities always completes this repair. Its total change is at most $h\eta$, giving a feasible lower value at least $Z_\eta-\lambda_0R-Kh\eta$. Source payoffs and path charges do not change. Their difference is at most $2Kh\eta=\varepsilon$.

Initialize every feasible source at layer zero and deficit zero, and process edges between consecutive layers. Arbitrary holes in the finite support of a preceding row are permitted by Lemma~\ref{lem:holes56}, yielding the stated concave convolution count. Direct enumeration of both deficits proves the nonconcave count. At $R=0$ there is one resource state and the method is exact. Numerical grid size, not its logarithm, appears in both bounds; rational bit complexity additionally depends on the encoding and exact evaluation of the supplied kernels.
\end{proof}

For renewal design, $H_u=\bar b-p(u)$ and $B_u=B-u$ at a feasible anchor, for which $p(u)=u$. The common deficit is therefore $\bar b-B$, and $h=m$ including the terminal edge. Prefix eligibility and concave rewards are needed to derive and implement the renewal arc functions; they are not needed for the recognition identity. A separate class of examples is modular resource allocation along a precedence-compatible configuration path, with $C_{uv}=p(v)-p(u)$ and separable concave returns on installed modules. These examples instantiate the abstract theorem without interpreting vertices as renewal commands. They are structural examples, not field-calibrated applications. Multiple independent resource equalities require a multidimensional state and componentwise repair; the scalar complexity bound must not be carried over unchanged.
\paragraph{The path quantifier.} The preceding recognition theorem concerns all source-to-sink paths in the retained graph. Under an edge allowance it is sufficient, but need not be necessary. The exact bounded-path recognition, witness and oracle-cost statements are in Theorem~\ref{thm:bounded60}; its layered construction does not add a second edge-count factor to the approximation recurrence.
\section{The Limited-Menu Complexity Frontier}\label{sec:param59}
Let $q_b=|\{b_j:j=1,\ldots,k\}|$ denote the number of distinct expected caps; the subscript distinguishes this count from a quadratic curvature. For a fixed command budget, enumerating all feasible books is polynomial with an exponent depending on that budget. Under a common reward, Theorem~\ref{thm:capcount56} reduces the useful book size to $\min\{m,2q_b\}$, and its coefficient two is tight. The following result identifies why neither parameter generally removes the exponent dependence.

\subsection{Encoded decision problem and the meaning of the parameters}\label{sec:decision60}
For the complexity results, an instance is an ordered list of rational catalog levels and nonnegative opening charges, positive rational probabilities summing to one, rational expected caps and realization ceilings, a rational accepted promise and value threshold $z$, an integer command allowance $m\geq1$, and the rational quadratic reward and service coefficients of Section~\ref{sec:model52}. Each rational is represented by a signed binary numerator and a positive binary denominator in lowest terms. Sorting and validation are polynomial-time preprocessing; duplicate catalog entries are consolidated at their least opening charge. The decision language asks whether an original feasible policy using a nonempty book of at most $m$ commands has net value at least $z$. Invalid or infeasible inputs are no-instances. Denote the total encoding length by $L$ and the number of distinct reduced rational caps by $q_b$.

Theorem~\ref{thm:wone59} is a parameterized many-one lower bound for this language with parameter $m+q_b$, already in the common linear-reward, zero-service-cost subclass. Hence it also excludes an algorithm $f(m)L^{O(1)}$ or $f(q_b)L^{O(1)}$, unless FPT equals W[1]: either such algorithm would be fixed-parameter in the sum. It does not assert W[1] membership or completeness, an ETH lower bound on the exact exponent, strong NP-hardness, or a kernel lower bound. The fixed-cap reconstruction and exact-lattice algorithms address different restrictions. In particular, numerical denominators cannot be replaced by their bit lengths in a pseudo-polynomial complexity expression. The positive result in Section~\ref{sec:tariff60} instead removes denominator dependence from the operation count by controlling opening-fee exceptions.
\begin{theorem}[Parameterized hardness in menu size and cap count]\label{thm:wone59}
The rational finite-catalog decision problem is $\mathrm{W}[1]$-hard parameterized by $m+q_b$. Hardness holds with the fixed common reward $f(x)=x$, zero service costs, uniform strictly positive probabilities, $\tau_j=b_j$, nonnegative rational opening charges, and $m=q_b$. Consequently, unless $\mathrm{FPT}=\mathrm{W}[1]$, no exact algorithm has running time $F(m,q_b)\mathcal L^c$ for an arbitrary computable $F$ and a constant $c$ independent of the parameters, where $\mathcal L$ is binary input length. The same conclusion holds for parameterization by either $m$ or $q_b$ alone.
\end{theorem}

The proof has two explicit stages. A group-sum encoding carries the consistency requirements of a multicolored clique. Dedicated baseline histories then make those group choices mandatory in the original renewal problem. Unlike a free augmentation under a nonbinding budget, every baseline must actually be installed in a threshold-achieving book.

\begin{lemma}[Group-sum encoding]\label{lem:groupsum59}
From an instance of multicolored clique with $r$ color classes, one can construct $g=r+\binom r2$ nonempty groups of positive integers $A_1,\ldots,A_g$ and a positive integer $W$, of total binary length polynomial in the graph description, such that a selection of at most one integer from each group sums to $W$ if and only if the graph has a multicolored clique. Every such sum uses exactly one integer from every group.
\end{lemma}
\begin{proof}
Empty color classes or a color pair with no edge give an immediate no-instance. Otherwise number vertices by distinct codes in $\{1,\ldots,n\}$. Create one group for each color and one for each unordered color pair. Use one digit for every group and two incidence digits for every color pair, in base $H>(2g+1)(n+1)$. The target has digit one at every group position and digit $n$ at every incidence position.

A vertex of color $i$ supplies digit one at its color-group position and its code at every incidence position involving that color. An edge joining codes $u$ and $v$ supplies digit one at its pair-group position and digits $n-u$ and $n-v$ at the two corresponding incidence positions. All unspecified digits are zero. Each encoded integer is positive because of its group digit. A selection of at most one integer per group has no carries: its digit sums are at most $gn<H$. Equality at the group digits forces every group to be represented. Equality at each incidence digit forces the chosen edge endpoint to equal the chosen vertex. Thus the selection is precisely a clique and its edges. Conversely such a clique supplies the required sum. There are $O(r^2)$ digits, each using $O(\log(rn))$ bits, and at most one encoded integer per vertex or edge. The construction is a parameterized reduction. Multicolored clique is $\mathrm{W}[1]$-hard: from ordinary $r$-clique make $r$ color copies of every vertex and connect distinct copies exactly when their original vertices are adjacent; clique hardness follows by complementing the independent-set completeness result of \citet{DowneyFellows1995b}.
\end{proof}

\begin{proof}[Proof of Theorem~\ref{thm:wone59}]
Apply Lemma~\ref{lem:groupsum59}. A target outside $[0,\sum_i\max A_i]$ is an immediate no-instance and can be mapped to the fixed feasible renewal instance with one zero command, cap and promise zero, and threshold one. Otherwise set
\[
 M=\max_i\max A_i,\qquad T=2(g+1)M,\qquad K=(2g+1)T,
 \qquad \ell_i=\frac{i}{g+1},\quad i=1,\ldots,g.
\]
There are $2g+1$ equally likely histories: one with cap and ceiling zero, and, for each group $i$, two with caps and ceilings $\ell_i$ and $\ell_i+\max A_i/T$. The catalog contains zero, every baseline $\ell_i$, and an optional command $\ell_i+w/T$ for each distinct $w\in A_i$. All points lie in $[0,1]$. Options within group $i$ lie strictly between its baseline and the next group's baseline; the last option is strictly below one. Equal integers within a group can be merged without changing the sum problem. Charge zero for zero and the baselines, and charge $w/(2K)$ for the option representing $w$. Set
\begin{equation}\label{eq:budgetreduction59}
 m=2g+1,\quad D_0=\frac{2}{2g+1}\sum_i\ell_i,\quad
 B=D_0+\frac{W}{K},\quad \zeta=D_0+\frac{W}{2K}.
\end{equation}
The $2g+1$ caps are distinct, so $q_b=m$. The promise is feasible because
$\bar b=D_0+\sum_i\max A_i/K\geq B$ and zero is in the catalog.

Any feasible book must contain zero, since the zero-cap history has positive probability. Write $S(c)$ for the sum of the integers represented by its installed optional commands. With linear reward and free service, its exact gross value is
\[
 \min\{B,D(c)\},\qquad D(c)=\sum_j\pi_j\max(c\cap[0,b_j]).
\]
If $B\leq D(c)$, lotteries between zero and each last eligible command attain the required aggregate terminal mean. If $B>D(c)$, use those last commands surely and fill the remaining promise with deterministic service; the available service capacity is $\bar b-D(c)$. These constructions obey the original realized ceilings, not only the expected caps.

For an upper bound only, augment the book by all baselines, ignoring the augmented count. Its terminal capacity is
\[
 D(c)\leq D_0+\frac{1}{K}\sum_i\max\bigl(\{w:\ell_i+w/T\in c\}\cup\{0\}\bigr)
 \leq D_0+\frac{S(c)}{K}.
\]
The actual book, not the augmented comparison, pays charge $S(c)/(2K)$. Hence
\begin{equation}\label{eq:keyparam59}
 V_c(B)\leq \min\{B,D_0+S(c)/K\}-S(c)/(2K)
       =\zeta-\frac{|S(c)-W|}{2K}.
\end{equation}
Suppose its value reaches $\zeta$. Then $S(c)=W$ and both capacity inequalities must be equalities. All option weights are positive, so the second equality allows at most one optional command per group. The first equality forces every baseline into the actual book: if $\ell_i$ is absent, the dedicated history with cap $\ell_i$ has a strictly smaller last command, and adding $\ell_i$ strictly increases its positive-weight contribution to $D(c)$. Zero and the $g$ baselines therefore consume $g+1$ slots. Lemma~\ref{lem:groupsum59} now forces one option from each group. These $2g+1$ installed commands are within the original budget and encode a clique.

Conversely, choose the $g$ options representing a clique together with zero and all baselines. This book uses exactly $m$ commands, has $D(c)=D_0+W/K=B$, and has net value $\zeta$. This proves equivalence. The rational construction has polynomial bit length, and $m+q_b=4g+2=O(r^2)$, proving the parameterized hardness statement.
\end{proof}

This theorem does not assert $\mathrm{W}[1]$ membership, a sharp exponent lower bound, strong NP-hardness, or hardness under uniform positive charges. It establishes a parameterized barrier for the original restricted model, rather than for a larger abstract path problem. Its fine rational scales also explain why it is consistent with the additive algorithm and the pseudo-polynomial lattice algorithm. Exact enumeration is XP, whereas an exponent independent of both menu size and cap count would collapse the stated parameterized classes. The full earlier SUBSET SUM construction and the tight cap-count proof are retained in the electronic companion.
\section{A Tariff-Sensitive Tractability Frontier}\label{sec:tariff60}
The hardness result uses command-specific opening charges. This section isolates a positive counterpart without bounding the command allowance, the number of distinct caps, or the numerical denominators. Throughout the section, the common reward is $f(x)=rx$ with $r>0$ and service is free, $h_j=0$. Expected caps, realization ceilings, probabilities and catalog spacing remain arbitrary subject to Section~\ref{sec:model52}.

\begin{lemma}[Terminal capacity of a book]\label{lem:capacity60}
For a feasible book $c=(u_1<\cdots<u_s)$, put
\[
 D(c)=\sum_j\pi_j\min\{b_j,d_j(c)\},\qquad
 d_j(c)=\max\{u\in c:u\leq\tau_j\}.
\]
Its exact gross value at the accepted promise $B$ is $r\min\{B,D(c)\}$. Moreover,
\begin{equation}\label{eq:gain60}
 D(c)=u_1+\sum_{i=1}^{s-1}G_{u_i u_{i+1}},\qquad
 G_{uv}=\sum_{j:\tau_j\geq v}\pi_j[\min\{b_j,v\}-u]_+.
\end{equation}
\end{lemma}
\begin{proof}
The aggregate terminal mean cannot exceed either $B$ or $D(c)$. If $B\leq D(c)$, distribute $B-u_1$ within the individual intervals $[u_1,\min\{b_j,d_j(c)\}]$, using their probability weights. Adjacent eligible commands implement each resulting terminal mean with zero service. If $B>D(c)$, first attain every terminal capacity. Histories with $b_j>d_j(c)$ then have a deterministic terminal command $d_j(c)$ and residual service capacity $b_j-d_j(c)$. These capacities total $\bar b-D(c)\geq B-D(c)$. Filling them meets the promise exactly and respects the realization ceiling because the final deterministic obligation is at most $b_j\leq\tau_j$. This attains the upper bound. For the identity, fix a history and telescope its eligible prefix after truncation at $b_j$. Every increment is precisely the summand in \eqref{eq:gain60}; the initial contribution is $u_1$ since $u_1\leq\min_j b_j$. Sum with weights $\pi_j$.
\end{proof}

Choose a standard opening fee $\rho_0\geq0$ and let $E=\{i:\rho_i\ne\rho_0\}$ and $d=|E|$. This parameter counts exceptional \emph{commands}, not distinct fee values. A most frequent catalog fee minimizes $d$ and is found without optimizing any allocation. For $J\subseteq E$, require exactly those exceptional commands in $J$ and forbid $E\setminus J$. Denote the allowed indices by $A_J$. An edge $i\to j$ with $i<j$ is admissible when $i,j\in A_J$ and no required index lies strictly between them. A starting index must precede or equal the first required index; a final index must follow or equal the last required index. These conventions are vacuous when $J$ is empty.

\begin{theorem}[Exact allocation with few tariff exceptions]\label{thm:tariff60}
Under the assumptions of this section, exact optimization and original-policy recovery require
\begin{equation}\label{eq:tariffcost60}
 O(kN^2+2^d mN^2)
\end{equation}
rational arithmetic operations, after replacing $m$ by $\min\{m,N\}$. The algorithm has bit complexity $2^d\operatorname{poly}(L)$ for binary input length $L$. It uses no discretization of the promise or capacities. In particular, a uniform nonnegative opening fee gives a polynomial-time algorithm with a strongly polynomial arithmetic count, for unrestricted $m$ and $q_b$.
\end{theorem}
\begin{proof}
Precompute all gains in \eqref{eq:gain60}. For each $J$, let $T^J_s(j)$ be the largest capacity of an admissible $s$-command prefix ending at $a_j$. Use $-\infty$ for absent states and initialize
\[
 T^J_1(j)=a_j
 \quad\text{if }j\in A_J,\quad a_j\leq\min\{B,\min_i b_i\},
 \quad j\leq\min J,
\]
where the last condition is omitted for $J=\varnothing$. The recurrence is
\begin{equation}\label{eq:tariffdp60}
 T^J_s(j)=\max_{i<j:\,i\to j\text{ admissible}}
       \{T^J_{s-1}(i)+G_{a_i a_j}\},\qquad 2\leq s\leq m.
\end{equation}
Accept only final indices $j\geq\max J$, with the condition again vacuous for an empty set. A valid start, edges that cannot skip a required index, and a valid end together force every member of $J$ to appear. Conversely, every book with exceptional set $J$ follows exactly such a path. For fixed $J$ and $s$, every accepted book pays the same charge
\[
 K^J_s=(s-|J|)\rho_0+\sum_{i\in J}\rho_i.
\]
By Lemma~\ref{lem:capacity60}, its value is nondecreasing in capacity. Keeping only the maximum capacity at each state therefore loses no optimal book. Maximize $r\min\{B,T^J_s(j)\}-K^J_s$ over accepted states and all subsets $J$. Predecessors recover an attaining book, and the lemma recovers a feasible original allocation. States with $s<|J|$ cannot reach a valid end and contribute nothing.

Computing the gains costs $O(kN^2)$. Prefix counts of required indices test edge admissibility in constant time, so each subset costs $O(mN^2)$. Capacity entries are sums of at most $N$ gains; equivalently, every attained entry is a sum of $k$ weighted truncated commands. A common product of the input denominators bounds their encoding length polynomially in $L$. Exact comparisons, the value calculation and the interval-filling recovery consequently have polynomial bit cost per operation. Neither a common denominator's numerical value nor a grid size enters \eqref{eq:tariffcost60}. This proves fixed-parameter tractability in $d$ and the uniform-fee conclusion.
\end{proof}

The distinction from Theorem~\ref{thm:wone59} is substantive. Its linear/free-service hard instances remain within this section's physical primitives but have command-specific fees that need not have few exceptions. Thus a cap-count and command-count lower bound coexists with an exception-count upper bound. Enumerating exceptions costs an exponential function of $d$ multiplying a polynomial whose exponent is independent of $d$; enumerating all books costs a parameter-dependent exponent in $N$. Uniform fees with nonlinear rewards or positive service costs are outside Theorem~\ref{thm:tariff60}, not instances silently solved by it.

\begin{corollary}[The uniform-fee decision frontier]\label{cor:fee60}
Fix the physical input and $B$. Let $D_s$ be the largest terminal capacity among feasible books with exactly $s$ commands, omitting infeasible sizes. Under a uniform fee $\rho$, the value is
\[
 V(\rho)=\max_{1\leq s\leq m}\{r\min(B,D_s)-s\rho\}.
\]
It is convex, nonincreasing and piecewise affine in $\rho$. The smallest optimal book size is nonincreasing in $\rho$, and every switching fee is an intersection of two displayed affine functions. The capacity tables compute the full fee frontier without rerunning the resource allocation problem.
\end{corollary}
\begin{proof}
The formula follows from the theorem. A maximum of decreasing affine functions is convex and nonincreasing. If $s_1,s_2$ are optimal at $\rho_1<\rho_2$, their two optimality inequalities imply $(s_1-s_2)(\rho_2-\rho_1)\geq0$. Pairwise intersections give all possible switches, with dominated intersections discarded. No diminishing-returns property in the command budget is assumed.
\end{proof}
\section{Merged-Origin Resource-Deficit Paths}\label{sec:deficit56}
The capacity potential does more than guarantee repair after rounding. It identifies a resource coordinate in which every feasible anchor has the same initial state and every feasible book has the same terminal state. This removes the extra anchor loop in the resource algorithm. The construction applies to the heterogeneous kernels of Theorem~\ref{thm:packing54}.

\subsection{An exact source-independent coordinate}
Write
\[
 p(u)=\sum_j\pi_j\min\{b_j,u\},\qquad p(\dagger)=\bar b,
 \qquad C_{uv}=T_{uv}+Q_{uv}=p(v)-p(u).
\]
For the terminal edge, $C_{u\dagger}=\bar b-p(u)$. Every feasible anchor $a\leq\min\{B,\min_jb_j\}$ satisfies $p(a)=a$. If edge $e$ is assigned resource $x_e$, define its unused capacity $w_e=C_e-x_e$ and its deficit payoff $\Psi_e(w)=\Phi_e(C_e-w)$. Opening charges are kept separate from $\Phi_e$ throughout. Let $R=\bar b-B$.

\begin{theorem}[Merged-origin identity]\label{thm:deficit56}
Under the common ordered catalog, prefix eligibility, and increasing concave history-specific rewards of Theorem~\ref{thm:packing54}, the original value equals
\begin{equation}\label{eq:deficit56}
 \max_{\substack{c:\,1\leq |c|\leq m\\a=\min c\leq\min(B,\min_j b_j)}}
 \left\{\sum_j\pi_j f_j(a)-\sum_{z\in c}\rho(z)
 +\max_{\substack{0\leq w_e\leq C_e\\\sum_{e\in P(c)}w_e=R}}
           \sum_{e\in P(c)}\Psi_e(w_e)\right\}.
\end{equation}
All anchors start at cumulative deficit zero and end at cumulative deficit $R$. Every feasible deficit vector on a chosen book yields a policy at the original promise, within the original command budget, with at least its displayed value. No anchor identifier is needed in a Bellman state once the current command and number of selected commands are known.
\end{theorem}
\begin{proof}
For a fixed book with anchor $a$, capacity conservation gives $\sum_e C_e=\bar b-a$. Hence
\[
 \sum_e x_e=B-a\quad\Longleftrightarrow\quad
 \sum_e(C_e-x_e)=\bar b-B=R.
\]
The transformation is a bijection between the two resource boxes intersected with their equalities. Apply the exact component-packing identity to obtain \eqref{eq:deficit56} and its recovery claim. Along a prefix ending at $u$, the cumulative deficit is
\[
 d=p(u)-a-\sum_{e\text{ in the prefix}}x_e.
\]
At the anchor it is zero; traversing $u\to v$ adds $C_{uv}-x_{uv}$. Continuation payoffs and transitions depend only on $u$, the remaining command allowance, and $d$, not on the earlier anchor. A Bellman merge retains the best accumulated value among such prefixes. The recovered full path retains its own anchor through predecessor pointers. Partial paths need not themselves encode a canonical history policy: the full-path packing theorem supplies implementation after the equality is restored.
\end{proof}

The same argument gives an abstract reuse principle. On an acyclic graph, suppose every resource capacity has the form $C_{uv}=p(v)-p(u)\geq0$, all paths share a sink, and a source $a$ requires resource $B-p(a)$ with $p(a)\leq B$. Complementing edge resources makes the total unused capacity $p(\dagger)-B$ independent of the source. Any additive path objective and edge-count bound then admit the same source merge. Concavity is needed for the approximation and fast convolution below, not for this algebraic identity. The renewal-specific content is the exact derivation of those capacities and payoffs from individual expected and realized constraints; an arbitrary constrained path need not have this potential structure.

\subsection{Centered approximation and exact repair}
For rational quadratic primitives put
\[
 r_{\max}=\max_jr_j,\qquad G=\max_j\gamma_j,
 \qquad \lambda_0=(r_{\max}-G)/2,\qquad K=(r_{\max}+G)/2>0.
\]
Every supergradient of $\Phi_e$ lies between $-G$ and $r_{\max}$ on its interval. Consequently
$g_e(w)=\Psi_e(w)+\lambda_0w$ is concave and $K$-Lipschitz. Adding this common linear term changes every exactly feasible path value by the same constant $\lambda_0R$; it does not change the optimizer.

For $\varepsilon>0$ and $R>0$, set
\begin{equation}\label{eq:mesh56}
 \eta=\frac{\varepsilon}{2Km},\qquad
 Q=\left\lfloor\frac{R}{\eta}\right\rfloor+1,\qquad
 \mathcal I=\{s\in\{0,\ldots,Q-1\}:R-m\eta\leq s\eta\leq R\}.
\end{equation}
Let $D_{\ell,v}(s)$ denote the best accumulated centered value at command $a_v$, with exactly $\ell$ selected commands and cumulative deficit $s\eta$. Unreachable entries are $-\infty$. All feasible anchors are initialized simultaneously:
\begin{align}
 D_{1,i}(0)&=\sum_j\pi_jf_j(a_i)-\rho_i,
       &&a_i\leq\min\{B,\min_jb_j\},\label{eq:init56}\\
 D_{\ell,v}(s)&=\max_{\substack{u<v,\ 0\leq t\leq s\\t\eta\leq C_{uv}}}
 \{D_{\ell-1,u}(s-t)+g_{uv}(t\eta)-\rho_v\},\label{eq:DP56}\\
 S_\eta&=\max_{\substack{1\leq\ell\leq m,\ u,\ s\in\mathcal I\\0\leq t\leq s,\ t\eta\leq C_{u\dagger}}}
 \{D_{\ell,u}(s-t)+g_{u\dagger}(t\eta)\}.\label{eq:terminal56}
\end{align}
The tail does not install a new command. A book with $s$ selected commands has $s$ resource edges including that tail.

\begin{theorem}[Centered deficit approximation]\label{thm:approx56}
Equations~\eqref{eq:init56}--\eqref{eq:terminal56} and exact fixed-book recovery produce an original-instance interval
\begin{equation}\label{eq:interval56}
 \underline V\leq V^*(B)\leq S_\eta-\lambda_0R+Km\eta,
 \qquad S_\eta-\lambda_0R+Km\eta-\underline V\leq\varepsilon.
\end{equation}
The guarantee allows heterogeneous rewards, arbitrary full-catalog eligibility classes, arbitrary nonnegative charges, and the original budget. For $R=0$ the same recurrence uses only deficit zero and is exact, including with nonzero service costs.
\end{theorem}
\begin{proof}
Round each deficit on an optimal continuous path down to a multiple of $\eta$. Its cumulative rounded deficit is no greater than $R$ and differs from $R$ by less than $m\eta$. Every prefix stays in the grid. Lipschitz continuity loses at most $Km\eta$ from the centered optimal value $V^*(B)+\lambda_0R$. Thus $S_\eta\geq V^*(B)+\lambda_0R-Km\eta$, proving the upper bound.

Conversely, recover a maximizing grid path. Increase its deficits, without exceeding their individual capacities, until their sum is $R$. This is always possible because the path's total capacity is $\bar b-a\geq\bar b-B=R$. The total increase is at most $m\eta$. The centered payoff decreases by at most $Km\eta$; after subtracting $\lambda_0R$, the uncentered path is worth at least $S_\eta-\lambda_0R-Km\eta$. Component packing gives an original feasible policy with no smaller value. Exact fixed-book allocation may improve it further, so the resulting feasible lower value satisfies the same bound. This proves \eqref{eq:interval56}. When $R=0$, nonnegative deficits must all be zero, so there is no rounding, repair, or approximation error.
\end{proof}

\subsection{Convolution with holes and bit complexity}
The merge can create unreachable gaps in a Bellman row. Concavity of the row itself is neither assumed nor generally true.

\begin{lemma}[Concave convolution with arbitrary finite support]\label{lem:holes56}
Let $A(z)\in\mathbb R\cup\{-\infty\}$ be arbitrary on $0,\ldots,Q-1$, and let $g(t)$ be a finite concave sequence on $0,\ldots,L$. The smallest maximizing input index in
\[
 M(s)=\max_{\max(0,s-L)\leq z\leq s}\{A(z)+g(s-z)\}
\]
is nondecreasing across rows having at least one finite feasible column. All values and predecessors can be computed in $O(Q\log(Q+1))$ comparisons and additions, including arbitrary holes in the support of $A$.
\end{lemma}
\begin{proof}
Delete columns with $A(z)=-\infty$ and rows whose feasible interval contains no remaining column. Concavity gives the anti-Monge inequality
\[
 [A(z)+g(s-z)]+[A(z')+g(s'-z')]
 \geq[A(z')+g(s-z')]+[A(z)+g(s'-z)]
\]
whenever $s<s'$, $z<z'$, and the four entries are feasible. If two maximizing indices crossed in the wrong order, their crossed rectangle would be feasible because $0\leq s-z\leq L$ is a staircase band. The inequality and the smallest-index tie rule would contradict that crossing. The maximizing indices are therefore monotone. Sorted finite columns permit feasible-row detection by interval searches. Divide-and-conquer row maximization scans a total $O(Q)$ candidate positions at each of $O(\log(Q+1))$ levels, with shared endpoint candidates included. Empty rows are discarded before recursion; they are not allowed to impose an artificial predecessor bound on another row.
\end{proof}

\begin{proposition}[Arithmetic, storage, and verification]\label{prop:cost56}
For rational quadratic inputs the centered deficit algorithm has arithmetic bound
\begin{equation}\label{eq:cost56}
 O\bigl(kN^2Q\log(k+1)+mN^2Q\log(Q+1)\bigr).
\end{equation}
Its cached implementation stores $O((N^2+mN)Q+kN^2)$ rational entries and has bit complexity polynomial in binary input length and numerical $Q$. The extra factor $N$ for separately enumerated anchors in \eqref{eq:fallback54} is absent. An independent checker can reconstruct all kernels and verify Bellman upper-potential inequalities in
$O(k^2N^2Q+mN^2Q^2)$ rational operations, without using the optimizer's monotone search.
\end{proposition}
\begin{proof}
There are $O(N^2)$ ordinary and terminal edges. Sorted weighted terminal segments and sorted service saturation thresholds give piecewise-linear terminal profiles and piecewise-quadratic service profiles. Their endpoints and cumulative values are rational. Direct evaluation by capped allocation already gives the first term in \eqref{eq:cost56}; caching endpoints and binary searching them can reduce its constant and history dependence. For each selected count and edge, Lemma~\ref{lem:holes56} gives the second term. The predecessor table has $mNQ$ entries. Caching every edge grid adds $N^2Q$ entries; retaining all component descriptions adds at most $kN^2$. Thus the storage claim includes the edge cache rather than counting Bellman states alone.

A kernel value involves finite sums and products of input rationals, a grid rational, and at most one capped-service water-level division. Sums of at most $k$ rational reciprocals have encoding length polynomial in the total input encoding. A path combines at most $m$ such values. Therefore all represented integers have bit length polynomial in the input length and $\log(Q+1)$; multiplying arithmetic counts by polynomial integer-arithmetic costs gives the stated bit bound. A common denominator may have exponentially large numerical value, which is not a polynomial-time claim in its binary encoding.

For independent verification, supplied Bellman values must dominate every feasible anchor and every predecessor-plus-edge transition. No transition into a claimed unreachable entry may be omitted. Every accepted terminal combination must be covered by the reported terminal bound. These inequalities inductively bound every grid path; an attaining path certifies the grid optimum used in the repair calculation. A separately checked original lottery supplies the lower value. Direct service-threshold testing costs at most quadratic time in histories per kernel value; explicitly examining every predecessor deficit gives the stated conservative checker bound. The checker uses neither the forward optimizer nor its divide-and-conquer routine.
\end{proof}

With $\delta=\varepsilon/K$ and $R\leq1$, substitution gives a budget-quadratic, catalog-quadratic term of order $m^2N^2\delta^{-1}\log(m/\delta+2)$ rather than its catalog-cubic predecessor. The guarantee is additive in signed net value. It is not polynomial in $\log(1/\varepsilon)$, and a successful optimization need not imply a cheap independent verification; both costs are measured separately.
\section{When a Price Bound Is Already an Exact Decision}\label{sec:pricebridge60}
A useful algorithmic distinction is not simply between an integer and a continuous method, but between a price bound that supports one implementable book and a bound that mixes books. Fix the original feasible book family $\mathcal C_B$. For $c\in\mathcal C_B$, let $F_c(t)$ denote its exact net value at aggregate target $t\in[c_1,\bar b]$. For a price $\lambda$, define
\[
 H_c(\lambda)=\max_t\{F_c(t)-\lambda t\},\qquad
 U(\lambda)=\lambda B+\max_{c\in\mathcal C_B}H_c(\lambda).
\]
A book is price-active if it attains the second maximum. The maximizing target set $I_c(\lambda)=\argmax_t\{F_c(t)-\lambda t\}$ is a closed interval by concavity. Its endpoints can be recovered by summing the smallest and largest individual maximizing targets. The exact price-path oracle in the companion computes $U(\lambda)$ without enumerating all books.

\begin{theorem}[Same-book support and a distance bound]\label{thm:pricebridge60}
For any price $\lambda$, $U(\lambda)=V^*(B)$ if and only if a price-active book $c$ has $B\in I_c(\lambda)$. For any price-active $c$, if its fixed-book value is $L$-Lipschitz, then
\begin{equation}\label{eq:pricedistance60}
 0\leq U(\lambda)-F_c(B)
 \leq (L+|\lambda|)\operatorname{dist}(B,I_c(\lambda)).
\end{equation}
For the quadratic primitives, the common bound in \eqref{eq:L52} is sufficient. With free service, common linear reward and zero opening fees, a single exact price certificate always exists: use $\lambda=r$ when $B\leq\max_c D(c)$, and $\lambda=0$ otherwise.
\end{theorem}
\begin{proof}
If $B\in I_c(\lambda)$ for an active book, then $U(\lambda)=F_c(B)\leq V^*(B)$, and weak duality gives equality. Conversely, take an optimal original book. Equality of the global bound with its value forces equality both in its individual price bound and in the maximization over books; it is active and its maximizing target interval contains $B$. For \eqref{eq:pricedistance60}, project $B$ onto $I_c(\lambda)$ at $T$. Then
$U(\lambda)-F_c(B)=F_c(T)-F_c(B)+\lambda(B-T)$, which gives the stated bound. A fixed-book allocation can be moved to a higher or lower aggregate target by filling or emptying residual individual target intervals. The total weighted change is exactly the aggregate change, so individual $L$-Lipschitz responses imply the same bound for $F_c$.

In the zero-fee linear case, $F_c(t)=r\min\{t,D(c)\}$ by Lemma~\ref{lem:capacity60}. At price $r$, every feasible book is active and its maximizing interval is $[c_1,D(c)]$. Choose a capacity-maximizing book when $B$ is below that capacity. At price zero, capacity-maximizing books are active on $[D(c),\bar b]$. This covers the remaining case.
\end{proof}

The interval must belong to \emph{one} active book. Merely placing $B$ between targets attained by different active books certifies a relaxation, not an original policy. Even a uniform positive fee can create this distinction. Take two equally likely histories with $b=\tau=(0,1)$, catalog $(0,1)$, $m=2$, common reward $f(x)=x$, free service, and fee $1/4$ for each installed command. At $B=1/8$, the one-command book has net value $-1/4$, whereas the two-command book has value $-3/8$. At $\lambda=1/2$, their price intercepts coincide at $-1/4$, producing the upper bound $-3/16$ and gap $1/16$. The tariff algorithm nonetheless solves this instance exactly. Thus its positive result is not merely a restatement of zero Lagrangian gap.

For method selection, this yields a concrete hierarchy. First check same-book support if a price bound is available. For common linear rewards and free service, a declared exception-count guard admits the exact tariff algorithm; exceeding that guard is an applicability decision, not a complexity guarantee. Outside this regime, price branching and the deficit approximation retain their original roles. A returned price interval supplies an instance-specific certificate. The deficit theorem supplies an a priori additive accuracy bound, but pays for a numerical resource state and, in the independent checker, all predecessor transitions. We compare these components separately rather than report an unimplemented hybrid as an experimental method.
\section{Selecting Reward Prototypes}\label{sec:grouping59}
The oscillation guarantee becomes an optimization tool only after groups and prototypes are chosen. For catalog reward vectors define
\[
 d(f,h)=\max_{a\in\mathcal A}[f(a)-h(a)]-\min_{a\in\mathcal A}[f(a)-h(a)].
\]
This is a pseudometric: symmetry follows by negating the difference, and the triangle inequality follows from subadditivity of maximum minus minimum. It is a metric after identifying vectors differing by a constant. Only histories with the same expected cap may share a prototype; ceilings, costs, probabilities, and promise remain unchanged.

For prototypes restricted to observed rewards, the minimum loss certificate with at most $G$ groups is the cap-restricted weighted medoid problem
\[
 \min\sum_{j,r:\,b_j=b_r}\pi_jd(f_j,f_r)x_{jr},\quad
 \sum_{r:b_r=b_j}x_{jr}=1,\quad 0\leq x_{jr}\leq z_r,\quad
 \sum_r z_r\leq G,\quad z_r\in\{0,1\}.
\]
For fixed selected centers a nearest-center assignment is optimal, so fractional assignment variables do not change the optimum. This is a standard facility-location model, not a new general clustering method. Its feasible upper bound supplies a valid loss certificate immediately, even if clustering is interrupted; a lower bound quantifies how much better this declared prototype family could be. Combined with Theorem~\ref{thm:robusttypes58}, a clustering solution of cost $\Delta$ and a prototype-policy solution of error $\xi$ give original regret at most $\Delta+\xi$ and an existence bound of $\min\{m,2G\}$ commands.

\begin{theorem}[Exact cap-preserving grouping frontier]\label{thm:grouping59}
Suppose that within cap class $b$, rewards are $f_j(x)=\alpha_j+r_jx-\kappa_b x^2/2$, with $r_j\geq\kappa_b\geq0$ and strict increase on the catalog. Among prototypes of this same class-specific curvature, the minimum oscillation loss for every group allowance through $G$ can be computed by weighted-median dynamic programming in
$O(k\log(k+1)+k^2G+q_bG^2)$ rational operations and $O(k^2+kG+q_bG)$ stored entries. This optimizes prototype selection, not the subsequent command-book search. Different expected caps are never merged.
\end{theorem}
\begin{proof}
Additive constants and common curvature cancel in each difference. Thus $d(f_j,f_r)=(a_N-a_1)|r_j-r_r|$. For a fixed cap class, sort the slopes. For any chosen ordered center slopes, nearest-center regions on a line can be chosen as contiguous intervals; weights multiply their distances but do not affect the nearest-center choice for an individual point. Conversely, each interval is optimally represented by a weighted median, which can be chosen among its observed slopes. The allowed median preserves increasing concavity. Therefore an optimum has contiguous groups with observed medians.

For sorted indices $1,\ldots,n_b$, let
$C_b(i,j)=(a_N-a_1)\min_r\sum_{v=i}^j\pi_v|r_v-r|$.
Prefix sums of weights and weighted slopes, with a median pointer that moves only forward as $j$ increases for fixed $i$, compute all interval costs in $O(n_b^2)$ operations. Define
\[
 F_b(t,j)=\min_{t-1\leq i<j}\{F_b(t-1,i)+C_b(i+1,j)\},
 \quad F_b(0,0)=0,
\]
with other infeasible states infinite. This yields every cap-class frontier in $O(Gn_b^2)$ operations. Convolve these frontiers over cap classes, allocating at least one group to each nonempty class and at most $G$ in total. This takes $O(q_bG^2)$ operations. Backtracking yields centers and assignments. The storage bound includes interval costs and backpointers. Summing over classes gives the result. With $G<q_b$ the declared cap-preserving grouping problem is infeasible; with $G\geq k$ its loss is zero.
\end{proof}

For arbitrary reward curves the medoid formulation remains valid, but the one-dimensional dynamic program is not asserted. No statistical estimation guarantee is implicit in this calculation. The tests compare the computed frontier with exhaustive medoid choices, and the resulting loss remains in the same reward-minus-cost units as the original objective.
\section{Computational Evidence and Method Selection}\label{sec:study59}
We distinguish mathematical identity checks, inherited evidence and new timed requests. All timings describe the recorded environments, not a claim of provider adoption or a distribution-free solver ranking. The electronic companion gives the nonrenewal decision mapping; the present experiment instantiates the service-contract model itself.

\subsection{Exact structural checks}
The new tariff algorithm agrees with exhaustive original-policy allocation on 172 small inputs, covering 6,302 feasible books. Five further exact comparisons use denominators with up to 513 bits. Forty-eight deliberately incomplete or corrupted certificates are rejected by a checker that never imports the tariff optimizer. The graph tests include 160 potential-generated graphs, 160 perturbed graphs with explicitly checked counterexample paths, and 480 length-bounded comparisons with exhaustive path enumeration. Four hundred rounding-and-repair cases preserve exact capacity bounds and equality, and 400 concave convolutions with arbitrary support holes agree with direct enumeration. The uniform-fee price-gap example is checked with exact fractions. These finite regressions support implementation correctness, not a proof by experiment. The full R59 hardness and prototype checks remain unchanged.

\subsection{What the inherited panel actually establishes}
The R59 panel has 72 tolerance-tagged cases but only 56 distinct model hashes: the baseline and two accuracy cells reuse the same eight specifications. It contains 288 method requests, not 288 independent instances. In the recorded three-second allowance, price search attains a rationally certified target on 71 of 72 cases, deficit on 5, and enumeration on 12. SCIP attains 72 numerical-bound targets; only its reconstructed lower policies have exact independent checks. Price search branches on seven cases, with maximum recorded depth 16. Its median end-to-end time is 0.339 seconds, compared with 1.852 for deficit, 1.827 for enumeration, and 0.527 for SCIP. These original measurements are retained, not rerun under a new environment label.

The deficit result is unfavorable as a general-purpose short-budget solver. Its 55 completed independent checks and 17 total deadlines do not become target attainment merely because a feasible fallback is available. Its numerical grid and the independent checker's predecessor enumeration explain a resource demand absent from a small price tree. Theorem~\ref{thm:pricebridge60} identifies a sufficient and necessary root-support condition; shallow observed search is not an empirical bound on arbitrary instances. The retained screening panel is likewise negative: total decision time improved in none of 18 matched pairs. Its safe-removal theorem remains useful, but the reported data do not justify enabling screening by default.

\subsection{Prospective standard-tariff and exception-tariff families}
Interpret a catalog level as a terminal service credit, an expected cap as the approved obligation for a history, and a realization ceiling as its maximum admissible realized obligation. The standard family has a common installation tariff, linear credit value and free preliminary service. The exceptions family changes exactly two catalog fees while retaining those physical primitives. The curved family has concave quadratic rewards and positive quadratic preliminary-service costs, so it deliberately lies outside the tariff theorem. These are stated model-derived interventions, not customer data or a borrowed external benchmark.

There are 30 seeds per family. Histories take sizes 8, 16 and 32; catalogs have 9, 17 and 25 levels; command allowances take values two, four and six. Caps lie on a thirty-second grid, ceilings add a nonnegative increment of at most four grid steps, and positive integer weights are normalized into probabilities. Promises take one quarter, one half or three quarters of aggregate capacity. Additive targets are one twentieth, one hundredth and one five-hundredth in the model's net-value units. The exact joint generator, including all interactions and fee values, is frozen before execution; these descriptive factors are not an orthogonal statistical design.

Every specification is run at total allowances of 0.5, 1.5 and 3 seconds. The internal optimization cutoff is 60 percent of each total, leaving the remainder for serialization and independent checking. Requests run sequentially in isolated processes; method order is shuffled deterministically within each case. Price, deficit, enumeration and the exact-formulation SCIP comparator run on every case. Tariff runs only in its two applicable families. No inapplicable case is counted as a tariff failure or as a success. The deficit state guard remains 600,000. All sources and inputs are committed before timing, and a completed request is never overwritten on restart.


% BEGIN retained/current source: revisions/or-r60-tariff-frontier-20260927/generated/new_outcomes.tex
The prospectively frozen study has 90 distinct input hashes, 270 budget-tagged cases and 1260 method requests. Tables~\ref{tab:newrational60} and~\ref{tab:newnumerical60} separate rational certification from numerical-bound evidence. The maximum uncompressed and compressed proof sizes are 339,597 and 55,483 bytes. All request times have median 0.347 seconds and 95th percentile 1.845 seconds, including deadlines. Standard: price 29/30, deficit 12/30, enumeration 19/30, SCIP 30/30, tariff 30/30 targets at three seconds. Exceptions: price 29/30, deficit 12/30, enumeration 19/30, SCIP 30/30, tariff 30/30 targets at three seconds. Curved: price 28/30, deficit 10/30, enumeration 18/30, SCIP 27/30 targets at three seconds. 

% END source


The full row-level file and family--budget--method summaries report the median, 90th and 95th percentiles, and maximum of total, optimization, serialization and checking time, alongside compressed and uncompressed proof bytes. Missing checking phases are not assigned zero cost: their denominator is recorded separately. Exact cap counts, joint types, weight-denominator bit lengths and predicted deficit-grid sizes accompany every input. Empirical attainment curves summarize each evidence class separately. Because the same seeds and specifications recur across interventions and budgets, these observations are dependent; the paper does not attach independent-sample confidence intervals to the request totals.

\subsection{Interpretation and decision scope}
The new tariff method tests the theorem's capacity compression, not a changed promise or a relaxed book budget. Its certificate checks every exception subset and capacity recurrence before checking an original policy. The SCIP formulation is algebraically exact, but its optimizer is numerical: we take the selected book and reoptimize its allocation in rational arithmetic. This reconstructed feasible policy may differ from the solver's raw continuous solution. Its exact lower value does not make the numerical upper bound exact. Accordingly, Tables~\ref{tab:newrational60} and~\ref{tab:newnumerical60} use separate primary evidence panels.

The results support regime-dependent method selection. A tariff applicability test, a one-book price certificate, and a resource approximation guarantee answer different questions. The new study does not establish that an automatic hybrid outperforms every component, since such a hybrid was not executed. Nor does a short-budget result measure the asymptotic benefit of longer branch-and-bound tuning. The exact theorem and the fully retained failures remain the basis for deciding which claims extend beyond the tested generators.
\section{Conclusion}
Joint allocation and implementation has a tariff-sensitive frontier. An exact selected-boundary representation preserves individual expected and realized constraints. In the common linear, free-service subclass, arbitrary opening fees retain parameterized hardness in command allowance and cap count; a bounded number of tariff exceptions instead yields fixed-parameter exact optimization, and uniform fees remove the resource grid entirely. The associated fee frontier translates that tractability into an installation decision.

For the broader concave model, capacity conservation and same-path repair retain the original-feasible additive guarantee. Length-aware recognition identifies the relevant path class rather than impose an unnecessary full-graph condition. Same-book price support explains exact certificates without confusing mixtures of books with implementable decisions. The computational evidence treats these as complementary regimes and preserves the deficit and screening methods' unfavorable observations.

Further distinctions remain mathematically consequential. Uniform positive fees with nonlinear rewards or costly service are not settled by the linear/free-service theorem. W[1] membership, strong-hardness and kernel questions require additional reductions or algorithms. Multiple independent resource promises require a multidimensional state. These boundaries do not remove the exact representation or the new positive regime; they specify the next structural questions beyond the proved frontier.

\paragraph{Code and data.} The accompanying current-submission archive contains standalone article and companion sources, the algorithms and independent checkers, prospectively frozen inputs, every completed or interrupted request, and reproduction instructions. A separate retained-reader directory preserves the predecessor manuscripts and their evidence. No historical timing is relabeled as a new execution.
\clearpage\phantomsection\label{refs-start}
\begin{thebibliography}{99}\raggedright
\bibitem[Aggarwal et al.(1987)Aggarwal, Klawe, Moran, Shor, and Wilber]{Aggarwal1987}
Aggarwal A, Klawe MM, Moran S, Shor P, Wilber R (1987) Geometric applications of a matrix-searching algorithm. \emph{Algorithmica} 2:195--208. doi:10.1007/BF01840359.
\bibitem[Ahmadi et al.(2025)Ahmadi, Raith, and Jalili]{Ahmadi2025}
Ahmadi S, Raith A, Jalili M (2025) A fast and simple algorithm for the resource constrained shortest path problem. \emph{33rd Annual European Symposium on Algorithms}, LIPIcs 351:97:1--97:15. doi:10.4230/LIPIcs.ESA.2025.97.
\bibitem[Arya et al.(2004)Arya, Garg, Khandekar, Meyerson, Munagala, and Pandit]{Arya2004}
Arya V, Garg N, Khandekar R, Meyerson A, Munagala K, Pandit V (2004) Local search heuristics for $k$-median and facility location problems. \emph{SIAM Journal on Computing} 33(3):544--562. doi:10.1137/S0097539702416402.
\bibitem[Babaioff et al.(2022)Babaioff, Gonczarowski, and Nisan]{Babaioff2016}
Babaioff M, Gonczarowski YA, Nisan N (2022) The menu-size complexity of revenue approximation. \emph{Games and Economic Behavior} 134:281--307. doi:10.1016/j.geb.2021.03.001.
\bibitem[Beasley and Christofides(1989)]{BeasleyChristofides1989}
Beasley JE, Christofides N (1989) An algorithm for the resource constrained shortest path problem. \emph{Networks} 19(4):379--394. doi:10.1002/net.3230190402.
\bibitem[Bolusani et al.(2024)]{Bolusani2024}
Bolusani S, et al. (2024) The SCIP Optimization Suite 9.0. ZIB Report 24-02, Zuse Institute Berlin. The execution manifest records the version actually run; this citation documents the solver framework.
\bibitem[Di Puglia Pugliese and Guerriero(2013)]{PuglieseGuerriero2013}
Di Puglia Pugliese L, Guerriero F (2013) A survey of resource constrained shortest path problems: Exact solution approaches. \emph{Networks} 62(3):183--200. doi:10.1002/net.21511.
\bibitem[Downey and Fellows(1995)]{DowneyFellows1995b}
Downey RG, Fellows MR (1995) Fixed-parameter tractability and completeness II: On completeness for W[1]. \emph{Theoretical Computer Science} 141(1--2):109--131. doi:10.1016/0304-3975(94)00097-3.
\bibitem[Fisher(1981)]{Fisher1981}
Fisher ML (1981) The Lagrangian relaxation method for solving integer programming problems. \emph{Management Science} 27(1):1--18. doi:10.1287/mnsc.27.1.1.
\bibitem[Gurski et al.(2019)Gurski, Rehs, and Rethmann]{Gurski2019}
Gurski F, Rehs C, Rethmann J (2019) Knapsack problems: A parameterized point of view. \emph{Theoretical Computer Science} 775:93--108. doi:10.1016/j.tcs.2018.12.019.
\bibitem[Handler and Zang(1980)]{HandlerZang1980}
Handler GY, Zang I (1980) A dual algorithm for the constrained shortest path problem. \emph{Networks} 10(4):293--309. doi:10.1002/net.3230100403.
\bibitem[Hart and Nisan(2019)]{HartNisan2013}
Hart S, Nisan N (2019) Selling multiple correlated goods: Revenue maximization and menu-size complexity. \emph{Journal of Economic Theory} 183:991--1029. doi:10.1016/j.jet.2019.07.006.
\bibitem[Hassin(1992)]{Hassin1992}
Hassin R (1992) Approximation schemes for the restricted shortest path problem. \emph{Mathematics of Operations Research} 17(1):36--42. doi:10.1287/moor.17.1.36.
\bibitem[Karp(1972)]{Karp1972}
Karp RM (1972) Reducibility among combinatorial problems. Miller RE, Thatcher JW, eds. \emph{Complexity of Computer Computations} (Plenum Press, New York), 85--103. doi:10.1007/978-1-4684-2001-2\_9.
\bibitem[Lorenz and Raz(2001)]{LorenzRaz2001}
Lorenz DH, Raz D (2001) A simple efficient approximation scheme for the restricted shortest path problem. \emph{Operations Research Letters} 28(5):213--219. doi:10.1016/S0167-6377(01)00069-4.
\bibitem[Maher et al.(2016)Maher, Miltenberger, Pedroso, Rehfeldt, Schwarz, and Serrano]{Maher2016}
Maher S, Miltenberger M, Pedroso JP, Rehfeldt D, Schwarz R, Serrano F (2016) PySCIPOpt: Mathematical programming in Python with the SCIP optimization suite. \emph{Mathematical Software---ICMS 2016}, 301--307. doi:10.1007/978-3-319-42432-3\_37.
\bibitem[Mohajerin Esfahani and Kuhn(2018)]{EsfahaniKuhn2018}
Mohajerin Esfahani P, Kuhn D (2018) Data-driven distributionally robust optimization using the Wasserstein metric: Performance guarantees and tractable reformulations. \emph{Mathematical Programming} 171:115--166. doi:10.1007/s10107-017-1172-1.
\bibitem[Pag\`es and Wilbertz(2012)]{PagesWilbertz2012}
Pag\`es G, Wilbertz B (2012) Intrinsic stationarity for vector quantization: Foundation of dual quantization. \emph{SIAM Journal on Numerical Analysis} 50(2):747--780. doi:10.1137/110827041.
\bibitem[Patriksson(2008)]{Patriksson2008}
Patriksson M (2008) A survey on the continuous nonlinear resource allocation problem. \emph{European Journal of Operational Research} 185(1):1--46. doi:10.1016/j.ejor.2006.12.006.
\bibitem[Rusmevichientong et al.(2010)Rusmevichientong, Shen, and Shmoys]{Rusmevichientong2010}
Rusmevichientong P, Shen ZJM, Shmoys DB (2010) Dynamic assortment optimization with a multinomial logit choice model and capacity constraint. \emph{Operations Research} 58(6):1666--1680. doi:10.1287/opre.1100.0866.
\bibitem[Sadykov et al.(2021)Sadykov, Uchoa, and Pessoa]{Sadykov2021}
Sadykov R, Uchoa E, Pessoa A (2021) A bucket graph-based labeling algorithm with application to vehicle routing. \emph{Transportation Science} 55(1):4--28. doi:10.1287/trsc.2020.0985.
\bibitem[Virtanen et al.(2020)Virtanen, Gommers, Oliphant, et al.]{Virtanen2020}
Virtanen P, Gommers R, Oliphant TE, et al. (2020) SciPy 1.0: Fundamental algorithms for scientific computing in Python. \emph{Nature Methods} 17:261--272. doi:10.1038/s41592-019-0686-2.
\bibitem[Winkler(1988)]{Winkler1988}
Winkler G (1988) Extreme points of moment sets. \emph{Mathematics of Operations Research} 13(4):581--587. doi:10.1287/moor.13.4.581.
\bibitem[Wu(1991)]{Wu1991}
Wu X (1991) Optimal quantization by matrix searching. \emph{Journal of Algorithms} 12(4):663--673. doi:10.1016/0196-6774(91)90039-2.
\bibitem[Yang et al.(2025)Yang, Song, Wang, Yang, and Leus]{Yang2025}
Yang K, Song G, Wang R, Yang H, Leus R (2025) Perspective Benders decomposition with applications to fixed-charge nonlinear resource allocation. \emph{INFORMS Journal on Computing}, published online July 8. doi:10.1287/ijoc.2024.0984.
\bibitem[Zhang(2012)]{Zhang2012}
Zhang H (2012) Solving an infinite horizon adverse selection model through finite policy graphs. \emph{Operations Research} 60(4):850--864. doi:10.1287/opre.1120.1056.


\end{thebibliography}
\label{refs-end}\clearpage
\begin{table}[p]\centering
\caption{Antecedents, retained results, and the tariff-sensitive frontier}\label{tab:frontier60}
\begin{tabular}{p{0.24\textwidth}p{0.31\textwidth}p{0.36\textwidth}}\toprule
Problem or method & Relevant antecedent & Result and boundary in this paper\\\midrule
Resource-constrained paths & Lagrangian bounds, labeling and approximation; Handler and Zang (1980), Hassin (1992), Sadykov et al. (2021), Ahmadi et al. (2025) & Original-model reduction with endogenous arc allocations, exact capacity conservation and same-path equality repair. No generic new labeling method.\\[0.07in]
Fixed-charge nonlinear allocation & Perspective reformulation and Benders cuts; Yang et al. (2025) & Ordered terminal-capacity compression; exact polynomial optimization at uniform fees and fixed-parameter optimization in exceptional commands.\\[0.07in]
Parameterized packing & Knapsack parameters and kernel questions; Gurski et al. (2019) & W[1]-hardness in command allowance plus cap count for the original feasible-policy language; no membership or kernel claim.\\[0.07in]
Finite support and prototypes & Moment-set geometry, dual quantization and facility location; Winkler (1988), Pag\`es and Wilbertz (2012), Arya et al. (2004) & Retained cap-count bound, oscillation transfer and cap-preserving grouping frontier; generic support and clustering facts are antecedents.\\[0.07in]
Price certificates & Standard weak duality and fixed-book concavity & Same-book support characterizes equality of a price bound; distance bounds and a positive-gap uniform-fee example distinguish it from tariff tractability.\\\bottomrule
\end{tabular}
\end{table}
\begin{table}[p]\centering\caption{Rationally certified target attainment by family and total allowance}\label{tab:newrational60}
\begin{tabular}{llrrr}\toprule Family & Method & 0.5 seconds & 1.5 seconds & 3 seconds\\\midrule
Standard & Price & 19/30 & 24/30 & 29/30\\
Standard & Tariff & 30/30 & 30/30 & 30/30\\
Standard & Deficit & 2/30 & 12/30 & 12/30\\
Standard & Enumeration & 10/30 & 17/30 & 19/30\\
Exceptions & Price & 19/30 & 24/30 & 29/30\\
Exceptions & Tariff & 30/30 & 30/30 & 30/30\\
Exceptions & Deficit & 2/30 & 12/30 & 12/30\\
Exceptions & Enumeration & 10/30 & 17/30 & 19/30\\
Curved & Price & 17/30 & 25/30 & 28/30\\
Curved & Tariff & -- & -- & --\\
Curved & Deficit & 0/30 & 5/30 & 10/30\\
Curved & Enumeration & 8/30 & 16/30 & 18/30\\
\bottomrule\end{tabular}\par\smallskip\begin{minipage}{\textwidth}\small Each denominator counts requests, not independent observations. Thirty specifications per family recur across budgets. Standard and exceptions have common linear rewards and free service; curved has quadratic rewards and positive service costs. Tariff is inapplicable, not failed, on curved instances. The numerical panel uses SCIP bounds with a reporting tolerance of $10^{-8}$; it is not a rational upper certificate.\end{minipage}\end{table}
\begin{table}[p]\centering\caption{Numerical-bound target attainment; independently checked lower policies only}\label{tab:newnumerical60}
\begin{tabular}{llrrr}\toprule Family & Method & 0.5 seconds & 1.5 seconds & 3 seconds\\\midrule
Standard & SCIP & 0/30 & 28/30 & 30/30\\
Exceptions & SCIP & 0/30 & 29/30 & 30/30\\
Curved & SCIP & 0/30 & 20/30 & 27/30\\
\bottomrule\end{tabular}\par\smallskip\begin{minipage}{\textwidth}\small Each denominator counts requests, not independent observations. Thirty specifications per family recur across budgets. Standard and exceptions have common linear rewards and free service; curved has quadratic rewards and positive service costs. Tariff is inapplicable, not failed, on curved instances. The numerical panel uses SCIP bounds with a reporting tolerance of $10^{-8}$; it is not a rational upper certificate.\end{minipage}\end{table}
\end{document}
